\section{Sparse pins on the torus}\label{sec:sparse-pins}

We now show that a sparse collection of prescribed slot points is
unlikely to lie on cycles avoiding the field pins.  This will supply
the geometric control of available slots used in
Section~\ref{sec:spatial-comparison}.  The induction follows the
pinned-boundary argument of \cite[Section~4]{BL}, itself developed from
\cite[Sections~5 and~7]{Ord}.  On the torus there is no pinned boundary:
a small uniform field supplies the killing, and cutting off flows to
infinity leaves an error that this killing must absorb.

\subsection{Soft pins and the ghost representation}

Rotate coordinates so that the field selects the up color.  If
\(q\in(0,1)\), \(w=1-q=e^{-b}\), its single-slot factor, after
removing a scalar, is
\begin{equation}\label{eq:soft-pin-factor}
 \diag(1,w)=wI+qP_{\mathrm{up}}.
\end{equation}
Thus we may independently select a pin on each real slot at the cut
with probability \(q\), and then weight the exchange picture by
\(2^{c}\), where \(c\) counts the cycles missing all pins.  Pinned
cycles have the up color; the other cycles receive independent fair
colors.  Expanding \eqref{eq:soft-pin-factor} proves that the resulting
spin law is the Gibbs law with exponent \(H_0+b\sum_xY_x\), after
rotating the up axis back to \(Y\).

For the induction we need a fixed pin set.  Add one stationary
\emph{ghost slot} for each real slot, pin each ghost at time
\(\beta/2\), and, just after the seam at time zero, independently
transpose each real slot with its ghost with probability \(q\).
The time-zero layer is taken after these switches.  A selected switch
inserts the corresponding ghost line into the real cycle; that cycle
then contains a ghost pin.  An unselected ghost remains a separate
pinned cycle and contributes weight one.  Consequently every switch
configuration has exactly the weight of the corresponding soft-pin
configuration.  This remains true with any finite set of deterministic
pins on real slot points.  Figure~\ref{fig:sparse-ghost} illustrates
the two cases.

\begin{figure}[ht]
 \centering
 \begin{tikzpicture}[x=1cm,y=1cm,>=Stealth,font=\small]
 \foreach \x/\lab in {0/{no switch},5.3/{switch}} {
  \node at (\x+1.3,3.25) {\lab};
  \draw[->,thick] (\x+.4,.4)--(\x+.4,2.5);
  \draw[->,thick,dashed] (\x+2.2,.4)--(\x+2.2,2.5);
  \node[above] at (\x+.4,2.5) {real};
  \node[above] at (\x+2.2,2.5) {ghost};
  \fill (\x+2.2,1.45) circle (2.5pt);
  \node[right] at (\x+2.2,1.45) {pin};
  \node[left] at (\x+.35,.4) {\(0^+\)};
  \node[left] at (\x+.35,2.3) {\(\beta^-\)};
 }
 \draw[thick] (.4,.4)--(.4,.05);
 \draw[thick,dashed] (2.2,.4)--(2.2,.05);
 \node at (1.3,-.3) {base weight \(w\)};
 \draw[thick] (5.7,.05)--(7.5,.4);
 \draw[thick,dashed] (7.5,.05)--(5.7,.4);
 \node at (6.6,-.3) {base weight \(q\)};
\end{tikzpicture}

 \caption{One real slot and its stationary ghost, shown on a cut-open
 time circle.  An absent switch leaves the ghost as a separate pinned
 cycle.  A present switch inserts its pinned line into the real cycle.
 In the unconditioned fresh real-hole traversal killed on entry into a ghost, passage through
 this layer survives with probability \(w=1-q\).  The rest of the
 real exchange picture is suppressed.}
 \label{fig:sparse-ghost}
\end{figure}

Let \(E\) denote all the ghost pins together with any finite set of
deterministic real-slot pins.  We write \(\PP_E\) for the normalized
picture law with density proportional to \(2^{c_E}\), and
\(\mathcal B_E\) for the slot-time points on cycles missing \(E\).
These cycles contain only real slots.  At coincident prescribed
operations, a slot point means the point just after the entire ordered
bundle, as in Section~\ref{sec:exchange}.  Distances
between sites are wrapped supremum distances, denoted \(d_L\); balls
are denoted by \(B_h(z)\).

Fix \(\alpha=1/100\) and a sufficiently large constant \(C_s\ge6\).
A finite collection of real slot points is \emph{sparse} if its site
coordinates, counted with multiplicity, satisfy
\begin{equation}\label{eq:torus-sparse}
 \#\{i:x_i\in B_h(z)\}\le C_s(1+h)^\alpha
 \qquad(z\in\Lambda_L,\ h\ge0).
\end{equation}
The same definition applies to lists of centers.  In particular,
different slots at the same site count separately.  The freedom to
choose \(C_s\) large will permit the finite collections of nearby
test points used in Section~\ref{sec:spatial-comparison}.

\begin{proposition}[Sparse-pin bound]\label{prop:sparse-pins}
Fix \(S\) and \(C_s\).  There is \(\beta_0<\infty\) such that, for
every fixed \(\beta\ge\beta_0\), there is \(L_*(\beta)<\infty\)
with the following property.  For every even \(L\ge L_*(\beta)\),
every \(q\) satisfying \(L^{-5/2}\le q<1\), every pin set \(E\)
as above, and every sparse set \(T\) of real slot points at one
deterministic time,
\begin{equation}\label{eq:sparse-pins}
 \PP_E(T\subseteq\mathcal B_E)\le p^{|T|},
 \qquad p=\beta^{-9/10}.
\end{equation}
The constants and thresholds are independent of the number, locations,
and times of the extra pins.
\end{proposition}

The proof is a finite downward induction on deterministic pins.
Its larger-pin hypothesis gives regular hole graphs; diffusion and a
flow-energy estimate then make next-pin returns rare enough to remove
the test pins by Lemma~\ref{lem:pin-inputs}.

If \(T\) is empty the assertion is immediate; if it meets \(E\),
the event is impossible.  We henceforth consider nonempty
\(T\cap E=\varnothing\), put \(P=E\cup T\), and expose the cycles
missing \(P\) under \(\PP_P\).  All other slot points are called
\emph{holes}.  Lemmas~\ref{lem:pin-inputs} and~\ref{lem:next-pin}
apply to this enlarged picture.  Here are the details concerning ghosts.

For any proposed set of exposed time-zero real labels, reveal its
trajectories, its seam images, and the switch outcomes encountered by
those trajectories.  Closure into cycles avoiding \(P\) requires every
such switch to be absent.  Separate Poisson clocks according to whether
they touch the currently exposed labels, as in the exposure
disintegration of Lemma~\ref{lem:pin-inputs}.  Conditional on the
revealed histories, the unused clocks are fresh on pairs of holes;
the remaining seam images are uniform completions within each site;
and switches on the other real--ghost pairs are still independent
Bernoulli switches.  When an exposed label jumps through a hole, that
hole is transported in the opposite direction.  A mark between two
exposed labels does not move a hole.  Requiring these to be all the
missing cycles conditions the fresh picture on every remaining cycle
meeting \(P\).  The exposed cycles contribute only their fixed factor
\(2^{c_P}\); therefore this description gives the conditional law,
including its all-hit conditioning.  That conditioning has positive
probability at almost every exposure.

Every ghost is a hole and is pinned.  Missing real cycles never use a
switch, so the number of real holes is constant in time.  The seam
does not change their number at any site.  In particular the
factorial-mark estimate in Lemma~\ref{lem:pin-inputs} remains valid:
inserting a real-edge transposition changes the number of missing
cycles by at most one, even in the presence of switches.

By Lemma~\ref{lem:next-pin}, conditional next-pin probabilities are
first-hit probabilities for the unconditioned fresh walk, with fresh
transitions on each repeated traversal.  Delete the other real-pin
stops and kill that walk on entering a ghost.  This only increases the
probability that its next pin belongs to \(T\): after entering a ghost,
the walk encounters its pin before its next opportunity to return to a
real slot.  The one-turn real-hole kernel at the time of \(T\) is
\begin{equation}\label{eq:ghost-mass}
 U=wK,\qquad U_+=U,\quad U_-=U^*,\quad Q=\tfrac12(U+U^*),
 \qquad E_y=\ip{y}{(I-Q)y}\ge q\norm{y}_2^2.
\end{equation}
Here and throughout this section matrices are column-forward.
The matrix \(K\) is doubly stochastic: its factors are symmetric
hole-walk heat kernels, forced transport bijections, and uniform
remaining seam completions.  Exactly one switch layer is crossed per
turn, and its real-to-real mean is \(wI\), which proves \(U=wK\).
All factors, their adjoints, and spatially killed versions contract
counting \(1\)-, \(2\)-, and supremum norms.  For example, a
nonnegative matrix whose row and column sums are at most one satisfies
\(\sum_i|(Av)_i|^2\le\sum_{i,j}A_{ij}|v_j|^2\le\norm v_2^2\).
This also proves the last inequality in \eqref{eq:ghost-mass}.

\subsection{The deterministic geometric input}

We state precisely the infinite-lattice results that will be transferred
to the torus.  An available-slot graph on \(\Z^3\) has between zero
and \(\ell\) slots at each site, all \(\ell\) slots outside a finite
set, and all pairs of available slots joined across every
nearest-neighbor site bond.  A site with no available slot is
\emph{blocked}.  Join blocked sites at supremum distance at most eight
and write \(\mathcal C(w)\) for the component of a blocked site \(w\).
Fix \(\nu=1/300\).  The constants \(M_*\) and \(A_0\) below depend
only on this lattice geometry and \(\ell\).

A graph is \emph{regular at \(z\) with radius \(A\ge A_0\)} if
\begin{enumerate}[label=(\roman*)]
 \item every ball \(B_{M_*}(v)\), with \(v\in B_{11A}(z)\),
 contains fewer than six blocked sites;
 \item each blocked \(w\notin B_{10A}(z)\) satisfies
 \(\max_{v\in\mathcal C(w)}|v-w|_\infty
       \le(1+|w-z|_\infty)^\nu\).
\end{enumerate}
The anchor itself need not have an available slot.
One may take \(M_*\ge\max\{48,\rho(3,44)+2\}\), where the
local-connectivity radius \(\rho(3,44)\) has the following property:
two unblocked sites at distance at most 44 can be joined inside
\(B_{\rho(3,44)}\) about the first site if that ball contains fewer
than six blocked sites.  The path has length at most
\((2\rho(3,44)+1)^3\).  This is
\cite[Lemma~5.5 and Equation~(5.13)]{Ord}; regularity is
\cite[Definition~3.2]{BL}, equivalently
\cite[Definition~5.3]{Ord}.  Take \(A_0\ge M_*+47\) large enough
for the following two conclusions.

\begin{lemma}[Lattice Nash inequality and unit flows]\label{lem:lattice-inputs}
In an available-slot graph regular at \(z\) with radius \(A\ge A_0\),
put
\[
 \mathcal D(v)=\tfrac12\sum_{\{i,j\}\text{ available edge}}
                              (v_i-v_j)^2.
\]
Every nonnegative \(v\) supported over \(B_A(z)\) satisfies
\begin{equation}\label{eq:local-nash}
 \norm v_2^{10/3}\le C\mathcal D(v)\norm v_1^{4/3}.
\end{equation}
For every available slot \(a\) at \(y\in B_A(z)\), there exists
an oriented edge flow \(F_a\), with divergence defined as outgoing
minus incoming flow, such that
\begin{equation}\label{eq:unit-flow}
 \Div F_a=\1_{\{a\}},\qquad
 |F_a(e)|\le C(1+\dist(e,y))^{-k},\qquad
 k=2-3\nu=199/100.
\end{equation}
Here \(\dist(e,y)\) is the minimum supremum distance to the two
endpoint sites.  The constants are uniform over the graph, \(A\),
\(z\), and the source.  Sites with any number of available slots
between one and \(\ell\) are permitted.
\end{lemma}

These are \cite[Proposition~3.3 and Lemma~3.4]{BL}; the underlying
geometric constructions are \cite[Propositions~5.6 and~6.2]{Ord}.
Their hypotheses are deterministic and impose no relation between
\(A\) and a time scale.  Only these stated conclusions of the
infinite-lattice results will be used here.

At a given exposure and physical time, make such a graph for each
torus anchor \(z\): lift the real holes in \(B_{\lfloor L/4\rfloor}(z)\)
to \(\Z^3\), and declare every slot outside this box available.
We say the slice is regular at \(z\) when this graph is regular.
We shall use
\begin{equation}\label{eq:sparse-localization-radius}
 A=\lceil C_{\mathrm{loc}}\beta\rceil,
 \qquad 11A+M_*<\lfloor L/4\rfloor,
\end{equation}
where \(C_{\mathrm{loc}}\) is fixed sufficiently large below.
Thus the Nash inequality for functions supported in \(B_A(z)\)
uses exactly the torus edges, with zero values beyond the support.

\subsection{From grid pins to regular slices}

Choose a finite circular grid \(\mathcal G\) containing zero, with
successive gaps at most \(p^{12}\).  For a set \(P\) of ghost pins
and real pins on this grid, let \(\mathsf I(P)\) assert
\eqref{eq:sparse-pins} under \(\PP_P\) for every sparse test at
one grid time.  In this subsection we assume \(\mathsf I(P)\).
This is the hypothesis for a larger pin set in the eventual downward
induction; no estimate for a smaller pin set is assumed.

At time \(t\), write \(\mathcal O_t\) for the sites without a real
hole.  At a grid time, simultaneous blocking of prescribed sparse
sites costs at most \(p^n\): choose one fixed slot at each site and
apply \(\mathsf I(P)\).  If any chosen slot is pinned, the event is
impossible.  The same conclusion can therefore be used without
checking disjointness separately.

\begin{lemma}[Blocking and component tails]\label{lem:torus-blocking}
Under \(\mathsf I(P)\), at any deterministic time \(t\),
\begin{align}
 \PP_P(A'\subseteq\mathcal O_t)&\le Cp^6
 &&\text{for six prescribed distinct sites},\label{eq:six-blocking}\\
 \PP_P(A'\subseteq\mathcal O_t)&\le(C\sqrt p)^{|A'|}
 &&\text{for a nonempty sparse set of sites}.\label{eq:sparse-blocking}
\end{align}
There are \(c,C,\lambda>0\), uniform for sufficiently large \(\beta\),
such that the blocked component formed with steps of supremum length
at most eight satisfies
\begin{equation}\label{eq:torus-component-tail}
 \PP_P\{z\in\mathcal O_t,\ 
       \max_{v\in\mathcal C_t(z)}d_L(z,v)\ge h\}
       \le C e^{-ch^\lambda}\qquad(h\ge1).
\end{equation}
All constants are independent of the grid and its size.
\end{lemma}

\begin{proof}
Let \(s\) be the grid point preceding \(t\), with \(s=t\) if
\(t\in\mathcal G\).  A site blocked at \(t\) either was blocked at
\(s\), or has an incident real-edge mark in \((s,t]\).  The seam
preserves hole counts at each site; the switch layer also preserves
the real-hole set of an exposure, as observed above.  The union of
the stars of six sites has bounded total intensity.  Factorial
domination in Lemma~\ref{lem:pin-inputs} bounds the probability of a
mark there by \(Cp^{12}\), giving \(p^6+Cp^{12}\le Cp^6\).

For \(n\) sparse sites, either at least \(\lceil n/2\rceil\)
were blocked at \(s\), costing at most
\(2^np^{\lceil n/2\rceil}\), or at least that many have incident
marks.  A mark touches at most two of the selected sites, so the
second alternative requires \(j=\lceil n/4\rceil\) distinct marks.
The union of the stars has intensity at most \(Cn\); hence its
probability is at most
\[
 \frac{(Cnp^{12})^j}{j!}\le(C'p^{12})^j,
\]
using \(j!\ge(j/e)^j\) and \(n/j\le4\).  Enlarging the constant
gives \eqref{eq:sparse-blocking}, without an independence assumption
between the two alternatives.

For the component tail, fix an integer \(a>100\) so large that
\(a^\alpha>4\), and set \(h_j=16a^j\).  A blocked path from
\(z\) to wrapped distance \(h_j\) yields two smaller crossings:
the first starts at \(z\), and the second at the first path site at
distance at least \(h_j/2\) from \(z\).  Truncate each at its
first reach of distance \(h_{j-1}\) from its own start.  The starts
are between \(h_j/2\) and \(h_j/2+8\) apart, and each truncated
crossing stays within radius \(h_{j-1}+8\).  Thus their mutual
distance is at least
\(h_j/2-2(h_{j-1}+8)\ge h_j/4\).  The remaining portion of the
original crossing has sufficient length for the second truncation,
since its endpoint is at least \(h_j/2-8\) from the second start.
These statements use only the triangle inequality for wrapped distance.

Recursion gives an ordered binary witness with \(2^j\) distinct
blocked leaves.  At every split of scale \(h_k\) the two descendant
leaf sets are at distance at least \(h_k/4\).  Therefore a ball of
radius less than \(h_k/8\) meets at most one branch at each split
of scale \(h_k\) or larger, and contains at most \(2^{k-1}\)
leaves.  For \(h_{k-1}/8\le h<h_k/8\), this is bounded by
\(C_s(1+h)^\alpha\), because \(a^\alpha>4\).  Below \(h_0/8\)
there is at most one leaf, and at or above \(h_j/8\) the same
comparison allows all \(2^j\) leaves.  Hence every such witness
is sparse, at every center and every radius.

At a split of scale \(h_k\), the second start has at most
\(Ch_k^3\) choices on the torus.  If \(W_j\) bounds the number
of witnesses rooted at a prescribed site, then
\[
 W_0=1,\quad W_j\le Ch_j^3W_{j-1}^2,\quad
 2^{-j}\log W_j\le\sum_{k\ge1}2^{-k}\log(Ch_k^3)<\infty.
\]
Thus \(W_j\le C_a^{2^j}\).  This count may ignore separation,
but the probability bound is applied only to the separated witnesses
just constructed.  Equation~\eqref{eq:sparse-blocking} bounds their
union by \((C_aC\sqrt p)^{2^j}\), which is at most
\(e^{-2\cdot2^j}\) for sufficiently large \(\beta\).
For \(h_j\le h<h_{j+1}\), use \(\lambda=\log_a2\) and
\(2^j\ge(h/(16a))^\lambda\).  Bounded \(h\) is absorbed by
\(C\); if \(h\) exceeds the torus diameter the event is empty.
This proves \eqref{eq:torus-component-tail}.
\end{proof}

\begin{lemma}[Regularity and localized diffusion]\label{lem:torus-diffusion}
Assume \(\mathsf I(P)\) and \eqref{eq:sparse-localization-radius}.
For every deterministic time and anchor, failure of regularity has
probability at most
\begin{equation}\label{eq:torus-irregularity}
 \Delta_A=CA^3p^6+C\exp(-cA^{\lambda\nu}).
\end{equation}
Call an exposure good at the anchor if its nonregular times have total
length at most \(\beta/100\).  Its probability of not being good is
at most \(100\Delta_A\).  On a good exposure, the one-turn kernel
localized by killing outside \(B_A(z)\) has every entry at most
\(C\beta^{-3/2}\), in either direction.

If a deterministic starting slot is a hole almost surely, localization
at its site loses at most \(Ce^{-c\beta}\) expected mass in any
forward or reversed traversal piece of length at most one turn.  This
last assertion does not require \(\mathsf I(P)\).
\end{lemma}

\begin{proof}
There are \(O(A^3)\) centers in \(B_{11A}(z)\), and a fixed
\(M_*\)-ball has only a bounded number of six-site subsets.
Equation~\eqref{eq:six-blocking} therefore gives \(CA^3p^6\)
for failure of local regularity.  For a remote blocked site at lifted
distance \(r>10A\), failure of the other condition costs at most
\(C\exp[-c(1+r)^{\lambda\nu}]\), by
\eqref{eq:torus-component-tail}.  Indeed all blocked sites of the
lift lie in its radius-\(\lfloor L/4\rfloor\) box, so coordinate
differences between them are at most \(L/2\) and equal their wrapped
distances.  A crossing in a lifted component is a crossing in the
torus component.  Summing the remote bound over shells, each containing
at most \(C(1+r)^2\) sites, gives the second term of
\eqref{eq:torus-irregularity}, with a smaller exponential rate.
Fubini and Markov give the asserted probability for a good exposure.

Let \(v(u)\) be a localized evolving nonnegative mass, initially of
mass at most one, and put \(H(u)=\norm{v(u)}_2^2\).  On smooth intervals,
with zero extension across the killing boundary,
\[
 H'=-2\mathcal D_t(v).
\]
The energy is the sum of \(\tfrac12(v_i-v_j)^2\) over unordered
real-hole edges, including edges from surviving to killed positions.
At a regular time, Lemma~\ref{lem:lattice-inputs} gives
\(H'\le-cH^{5/3}\).  At other times and at instantaneous
factors, the squared norm cannot increase.  Unless it has become zero,
\((H^{-2/3})'\ge c'\) throughout the regular times.
Each half turn has at least \(\beta/2-\beta/100\) regular time,
so its output has \(2\)-norm at most \(C\beta^{-3/4}\).
The same proof applies to its adjoint, since heat factors are
symmetric, transports reverse by bijections, and seam completions
reverse by uniform inverse completions.  Split a full turn at its
temporal midpoint and apply Cauchy--Schwarz to the two resulting
masses.  This proves the entry bound.

For localization, sample the complete picture under \(\PP_P\),
and independently supply fresh Poisson clocks for the hole walk,
accepted where both endpoints are holes.  Forced transports use
actual picture marks; seam and switch operations do not change sites.
An exit from \(B_A(z)\) needs at least \(A\) inter-site steps.
In a traversal piece of length at most \(\beta\), the physical
times run injectively through at most one period, whether forward,
backward, or across the seam.  Thus the actual marks used by these
steps are distinct, as are the fresh marks.  Ignore slot compatibility
and count site paths.  There are boundedly many choices of the next
site, slot edge, and actual-or-fresh clock at each step.  The factorial
bound of Lemma~\ref{lem:pin-inputs}, combined with the independent
fresh factorial measure, bounds the expected number of ordered lists
of \(j\) steps from the starting site by
\[
 \frac{(C\beta)^j}{j!}.
\]
The denominator is the volume factor for the ordered time simplex.
Taking \(j=A\) and \(C_{\mathrm{loc}}\) sufficiently large makes
this at most \(Ce^{-c\beta}\), using \(j!\ge(j/e)^j\).
This is an unconditional estimate under the correct exposure marginal;
no fresh-clock estimate under the all-hit conditioning is needed.
\end{proof}

\subsection{Cutting off the flows}

Localized diffusion controls a direct one-turn return.  A return after
many turns will instead be controlled by an energy inequality.  The
following finite-torus form of the sparse-source estimate isolates the
only error introduced by replacing a pinned boundary with ghost killing.

\begin{lemma}[Sparse-source pairing on the torus]\label{lem:torus-flow-pairing}
At a fixed slice let \((z_i)_{i=1}^{m_T}\) be a sparse list, including
repetitions.  Retain any subset of its members at which the lifted
graph is regular with radius \(A\).  For each retained member let
\(\mu_i\ge0\) be a mass at most one, supported on the real holes
over \(B_A(z_i)\).  If \(L\) satisfies
\eqref{eq:sparse-localization-radius} and is sufficiently large
compared with \(A\), then every real vector \(Y'\) on the holes
of that slice satisfies
\begin{equation}\label{eq:torus-flow-pairing}
 \left|\left\langle\sum_i\mu_i,Y'\right\rangle\right|
 \le (Cm_TA^\alpha)^{1/2}\sqrt{\mathcal D_t(Y')}
      +Cm_TL^{1/2-k}\norm{Y'}_2.
\end{equation}
\end{lemma}

\begin{proof}
For an individual source slot \(a\) in the support of \(\mu_i\),
take the flow of Lemma~\ref{lem:lattice-inputs} in the lift about
\(z_i\).  Multiply its value on each edge by a taper equal to one
within anchor radius \(L/10\), zero beyond radius \(L/5\), and
with Lipschitz constant \(C/L\).  For instance use a piecewise linear
function of the supremum distance of the edge midpoint.  The support
is strictly inside the lifted box, so this cut flow projects to a flow
\(\widetilde F_a\) on the torus real-hole graph.

Its divergence has the form \(\1_{\{a\}}+\varepsilon_a\).
To see the error bound, compare the taper at an incident edge with
its value at the vertex.  Their difference is at most \(C/L\);
the untapered divergence is zero except at \(a\), where the taper
is one.  In the transition annulus every edge is at distance comparable
with \(L\) from the source, because \(A\ll L\).  The bounded
degree and \eqref{eq:unit-flow} give
\[
 |\varepsilon_a(j)|\le CL^{-k-1},\qquad
 \norm{\varepsilon_a}_2\le CL^{1/2-k}.
\]
No discrepancy arises outside that annulus enlarged by one lattice
step, and there are at most
\(CL^3\) slots in it.

The projected flow retains the bound
\(C(1+d_L(e,x_a))^{-k}\).  For two source sites at wrapped
distance \(h\), the sum of the product of these envelopes is at most
\begin{equation}\label{eq:torus-flow-overlap}
 C(1+h)^{-(2k-3)}.
\end{equation}
For \(h\ge3\), inside the radius-\(h/3\) ball about either source,
one envelope is at most \(Ch^{-k}\) and the sum of the other is at
most \(Ch^{3-k}\).  Outside both balls, dyadic shells of radii at
least \(h/3\) contribute at most
\(C\sum_{r\ge0}(2^rh)^{3-2k}\).  Shell counts on the torus are
bounded by the corresponding lattice counts.  Since \(2k>3\),
this proves \eqref{eq:torus-flow-overlap}; bounded \(h\) follows
from square summability of the envelopes.

Sum the cut flows with the source masses as coefficients and denote
the result by \(F\).  For centers at distance at most \(4A\),
use the constant overlap bound; for more distant centers, every pair
of their source sites is at least half as far apart.  With
\(\sigma=2k-3=49/50>\alpha\), sparseness and dyadic annuli give
\begin{align*}
 \norm F_2^2
 &\le C\sum_i\left(\#\{j:d_L(z_i,z_j)\le4A\}
       +\sum_{j:d_L(z_i,z_j)>4A}(1+d_L(z_i,z_j))^{-\sigma}\right)\\
 &\le Cm_T\left(A^\alpha+
          \sum_{r\ge0}(2^rA)^{\alpha-\sigma}\right)
 \le Cm_TA^\alpha.
\end{align*}
This calculation uses only that each mass is at most one; it permits
both overlapping supports and repeated centers.  The summed divergence
error \(\varepsilon\) has
\(\norm\varepsilon_2\le Cm_TL^{1/2-k}\).  Finally,
\[
 \left\langle\sum_i\mu_i,Y'\right\rangle
 =\sum_eF(e)(Y'_{e^-}-Y'_{e^+})-\ip{\varepsilon}{Y'}.
\]
Cauchy--Schwarz, with
\(\sum_e(Y'_{e^-}-Y'_{e^+})^2=2\mathcal D_t(Y')\), proves the
claim.  The empty retained subset gives zero and is included.
\end{proof}

\subsection{Return energy and completion of the induction}

We now prove Proposition~\ref{prop:sparse-pins}.  Fix a grid and
argue by downward induction over its finite set of real slot points.
When every such point is pinned, \(\mathsf I(E)\) is immediate.
Assume the assertion for every strict superset of \(E\), and choose
a nonempty sparse grid test \(T\) disjoint from \(E\).  Then
\(P=E\cup T\) satisfies \(\mathsf I(P)\), so all the preceding
blocking and diffusion estimates hold under \(\PP_P\).

At an exposure under that law put \(m_T=|T|\) and
\(\zeta=\1_T\).  Let \(F=1\) on \(T\), and elsewhere let
\(F\) be the probability that the ghost-killed walk with kernel
\(U\) ever hits \(T\) at a later cut.  Explicitly, start with
\(F_0=\zeta\) and iterate
\(F_{n+1}=1\) on \(T\), \(F_{n+1}=U^*F_n\) off \(T\).
These vectors increase to \(F\), with \(0\le F\le1\) and
\(U^*F=F\) off \(T\).  Set \(g=F-\zeta\), so
\(0\le g\le1\) and \(g=0\) on \(T\).

Let \(r(E,T)\) be the mean, over \(\PP_P\) and the uniform
starting point in \(T\), of the indicator that its next pin is
in \(T\).  The pin inequality of Lemma~\ref{lem:pin-inputs} is
\begin{equation}\label{eq:sparse-pin-reduction}
 \PP_E(T\subseteq\mathcal B_E)\le(2r(E,T))^{m_T}.
\end{equation}
Conditionally on the exposure, deletion of the other real-pin stops
bounds the sum of these next-pin probabilities by
\(\sum_{i\in T}(U^*F)_i\).  Harmonicity off \(T\) gives
\begin{align}
 \sum_{i\in T}(U^*F)_i
 &=m_T-\ip F{(I-Q)F}\notag\\
 &=\ip\zeta{Q\zeta}+2\ip{Q\zeta}g-E_g.
 \label{eq:torus-hitting-energy}
\end{align}
Indeed the quadratic forms of \(U,U^*,Q\) agree on real vectors,
and \(\ip\zeta g=0\).  In particular no reversibility of \(U\)
has been assumed.

The first term is controlled by one-turn diffusion.  Start a unit mass
at each \(i\in T\) and localize at its site, in either direction.
Lemma~\ref{lem:torus-diffusion} bounds discarded mass in expectation
by \(Ce^{-c\beta}\).  On a good exposure there are at most
\(C_s(1+A)^\alpha\) target points of \(T\) in its localized
region, each receiving at most \(C\beta^{-3/2}\).  A nongood
exposure has probability at most \(100\Delta_A\) and still carries
total mass at most one.  Thus
\begin{equation}\label{eq:torus-short-return}
 \E_P\ip\zeta{Q\zeta}
 \le Cm_T\bigl(e^{-c\beta}+\Delta_A+A^\alpha\beta^{-3/2}\bigr).
\end{equation}

To absorb the linear term in \eqref{eq:torus-hitting-energy}, split
the traversal in direction \(d\in\{+,-\}\) at elapsed time
\(u\in I=[\beta/3,2\beta/3]\):
\[
 U_d=B_d(u)C_d(u),\qquad Y_d(u)=B_d(u)^*g.
\]
The factors use matching hole layers at physical time
\(t_d(u)=\theta+u\) or \(\theta-u\) modulo \(\beta\), where
\(\theta\) is the cut of \(T\).  Each instantaneous factor is
assigned to exactly one side; the finitely many exceptional split
times do not affect the integrals.  As \(u\) decreases, \(Y_d(u)\)
accumulates adjoint factors from the endpoint.  Its continuous norm
loss is twice its Dirichlet energy, and instantaneous factors contract.
Consequently
\begin{equation}\label{eq:torus-suffix-energy}
 \int_I\mathcal D_{t_d(u)}(Y_d(u))\dd u
 \le\tfrac12\bigl(\norm g_2^2-\norm{U_d^*g}_2^2\bigr)
 \le E_g.
\end{equation}
The second inequality follows from the exact identity
\[
 E_g-\tfrac12(\norm g_2^2-\norm{U_d^*g}_2^2)
       =\tfrac12\norm{g-U_d^*g}_2^2,
\]
because \(\ip g{U_d^*g}=\ip g{Qg}\).
Also \(\norm{Y_d(u)}_\infty\le1\) and
\(\norm{Y_d(u)}_2\le\norm g_2\).

For each \(i\in T\), propagate its unit mass by the starting
portion \(C_d(u)\), with localization at \(z_i=x_i\), and retain
it only when the slice is regular at \(z_i\).  Call the resulting
sum \(M_d(u)\), and put
\[
 J_d=|I|^{-1}\int_I\ip{M_d(u)}{Y_d(u)}\dd u.
\]
The original pairing
\(\ip{C_d(u)\zeta}{Y_d(u)}=\ip{U_d\zeta}g\) is constant in
\(u\).  Supremum contraction, localization, and the deterministic-time
bound \eqref{eq:torus-irregularity} show that
\begin{equation}\label{eq:torus-pairing-error}
 \E_P\left|\ip{U_d\zeta}g-J_d\right|
 \le Cm_T(e^{-c\beta}+\Delta_A).
\end{equation}
This uses Fubini and only regularity at the split, without requiring a
good exposure over the entire turn.

Apply Lemma~\ref{lem:torus-flow-pairing} to the retained masses.
Cauchy--Schwarz in \(u\), \eqref{eq:torus-suffix-energy}, and
\(E_g\ge q\norm g_2^2\) give the pointwise bound
\begin{equation}\label{eq:torus-linear-energy}
 |J_d|\le C\left((m_TA^\alpha/\beta)^{1/2}
           +m_TL^{1/2-k}/\sqrt q\right)\sqrt{E_g}.
\end{equation}
Since \(2\ip{Q\zeta}g=\ip{U_+\zeta}g+\ip{U_-\zeta}g\),
we complete the square in \(\sqrt{E_g}\) to obtain
\[
 J_++J_--E_g
 \le C\left(m_TA^\alpha/\beta+m_T^2L^{1-2k}/q\right).
\]
Combining this with \eqref{eq:torus-short-return} and
\eqref{eq:torus-pairing-error}, and dividing by \(m_T\), yields
\begin{equation}\label{eq:torus-return-bound}
 r(E,T)\le C\left(e^{-c\beta}+\Delta_A
     +A^\alpha\beta^{-3/2}+A^\alpha/\beta
     +m_TL^{1-2k}/q\right).
\end{equation}

Every constant so far is independent of the grid, its number of
points, and the induction pin set.  Sparseness gives
\(m_T\le C L^\alpha\); because \(q\ge L^{-5/2}\), the last term
of \eqref{eq:torus-return-bound} is at most
\[
 C L^{\alpha+7/2-2k}=C L^{-47/100}.
\]
The remaining terms, with \(A=\lceil C_{\mathrm{loc}}\beta\rceil\),
are \(o(p)\) as \(\beta\to\infty\): their polynomial orders are
\(\beta^{-12/5}\), \(\beta^{-149/100}\), and
\(\beta^{-99/100}\), besides stretched-exponential terms.
Choose \(\beta_0\) sufficiently large first.  For each fixed
\(\beta\ge\beta_0\), choose \(L_*(\beta)\) sufficiently large
for \eqref{eq:sparse-localization-radius} and for the last term to
make \(2r(E,T)\le p\).  Equation~\eqref{eq:sparse-pin-reduction}
then proves \(\mathsf I(E)\), completing the finite induction.

For arbitrary finite extra pins and a test at any deterministic time,
choose a grid containing their times and zero, and refine it to mesh
at most \(p^{12}\).  Uniformity in the grid gives the proposition
without a limiting procedure.  In particular the blocking,
component-tail, and regularity conclusions above now hold under every
allowed pin law.  All time integrals are measurable: the finite-volume
histories have finitely many jumps, their kernels are products of
measurable heat and instantaneous matrices, and the hitting function
is an increasing limit of finite-horizon probabilities.  No choice of
the auxiliary flows has to be measurable, since only their deterministic
pairing inequality is used.

\begin{corollary}[Positive spontaneous magnetization]\label{cor:positive-magnetization}
For every fixed \(S\) and all sufficiently large finite \(\beta\),
\[
 m_{S,\beta}\ge S(1-\beta^{-9/10})>0.
\]
\end{corollary}

\begin{proof}
The passage from pinned-cycle probabilities to the pressure slope uses
pressure increments, as in \cite[Section~7]{BL}.  Fix \(\beta\)
above the threshold of Proposition~\ref{prop:sparse-pins}, enlarged
so that \(\beta>1\), and put \(p=\beta^{-9/10}\).  For an exponent
field \(s>0\), set \(q(s)=1-e^{-s}\), and use the ghost picture
with no extra real pins.  A real slot is down with probability one
half of its probability of belonging to a cycle missing the ghost
pins: pinned cycles are up, and missing cycles are colored fairly.
The singleton case of Proposition~\ref{prop:sparse-pins} therefore
gives a down probability at most \(p/2\), whenever
\(L\ge L_*(\beta)\) and \(q(s)\ge L^{-5/2}\).

Define the finite-volume exponent pressure
\[
 \mathcal P_L(s)=V^{-1}\log\Tr
       \exp\!\left(H_0+s\sum_xY_x\right).
\]
Write \(\langle\cdot\rangle_s\) for the corresponding tilted Gibbs
expectation.
There are \(\ell V\) slots and \(\ell=2S\), so differentiation
and the color representation give
\[
 \mathcal P_L'(s)
   =V^{-1}\left\langle\sum_xY_x\right\rangle_s
   \ge S(1-p).
\]
For fixed \(0<a<b\), this holds simultaneously for all
\(s\in[a,b]\) once
\(L\ge\max\{L_*(\beta),q(a)^{-2/5}\}\).  Integration yields
\[
 \mathcal P_L(b)-\mathcal P_L(a)\ge S(1-p)(b-a).
\]
By rotation invariance and the pressure limit \eqref{eq:pressure},
\(\mathcal P_L(s)\to\mathcal P(s)=\beta p_{S,\beta}(s/\beta)\)
along even \(L\).  Pass to this limit first.  The limiting pressure
is \(S\)-Lipschitz, because
\(\norm{\sum_xY_x}\le SV\); letting \(a\downarrow0\) therefore
gives \(\mathcal P(b)-\mathcal P(0)\ge S(1-p)b\).
Finally divide by \(b\) and let \(b\downarrow0\).  Since
\(\mathcal P'(0+)=p_{S,\beta}'(0+)=m_{S,\beta}\), this proves
the claim with the volume limit taken before removal of the field.
\end{proof}
